\usepackage{ifthen}
\usepackage{xifthen}
\usepackage{epigraph}
\usepackage{stackengine}

\def\testdateparts#1{\dateparts#1\relax}
\def\dateparts#1 #2 #3 #4 #5\relax{
    \marginpar{\small\textsf{\mbox{#1 #2 #3 #5}}}
}


\makeatletter
\def\@lesson{}%
\newcommand{\lesson}[3]{
    \ifthenelse{\isempty{#2}}{%
        \ifthenelse{\isempty{#3}}{%%
            \addcontentsline{toc}{subsubsection}{Lesson #1}
            \def\@lesson{Lesson #1}%%
        }{%%
            \addcontentsline{toc}{subsubsection}{Lesson #1: #3}
            \def\@lesson{Lesson #1: #3}%%
        }%%
    }{%
        \ifthenelse{\isempty{#3}}{%%
            \addcontentsline{toc}{subsubsection}{Lesson #1}
            \def\@lesson{Lesson #1 \normalfont\small{\mbox{-- #2}}}%%
        }{%%
            \addcontentsline{toc}{subsubsection}{Lesson #1: #3}
            \def\@lesson{Lesson #1: #3 \normalfont\small{\mbox{-- #2}}}%%
        }%%
    }%

    \subsection*{\@lesson}\
}
\makeatletter

\renewcommand\date[1]{\marginpar{#1}}

\RequirePackage{subfiles}

\usepackage[utf8]{inputenc}
\usepackage[T1]{fontenc}
\usepackage{lmodern}

\usepackage[fixamsmath]{mathtools}
\usepackage[nointegrals]{wasysym}
\usepackage{physics}
\usepackage{adjustbox}

\usepackage{XCharter}
\usepackage{fourier}
\usepackage{amsmath,amssymb,amsthm}
\usepackage[cal=cm, bb=ams]{mathalpha}

\usepackage[margin=0.75in]{geometry}
\usepackage[dvipsnames, svgnames, table]{xcolor}
\usepackage{quiver}

\definecolor{PageBackground}{HTML}{1E1E1E}
\definecolor{PageText}{HTML}{E0E0E0}

\pagecolor{PageBackground}
\color{PageText}

\usepackage{graphicx}
\usepackage{tikz}
\usepackage{tikz-cd}
\usetikzlibrary{arrows.meta, cd}
\usepackage{pgffor}

\tikzset{every picture/.style={text=PageText}}

\usepackage{xparse}
\usepackage{enumitem}
\setlist[itemize]{itemsep=0.1pt, topsep=0pt, partopsep=0pt, label=\textcolor{PageText}{$\bullet$}}
\setlist[enumerate]{itemsep=0.1pt, topsep=0pt, partopsep=0pt, label=\textcolor{PageText}{\arabic*.}}
\usepackage{cancel}
\usepackage{hwemoji}

\DeclareUnicodeCharacter{FE0E}{}
\DeclareUnicodeCharacter{FE0F}{}

\usepackage[most]{tcolorbox}
\tcbuselibrary{skins}

\definecolor{AxiomBar}{HTML}{2EC4B6}
\definecolor{TheoremBar}{HTML}{FF6B6B}
\definecolor{PropositionBar}{HTML}{FF9F43}
\definecolor{LemmaBar}{HTML}{A66CFF}
\definecolor{CorollaryBar}{HTML}{38BDF8}
\definecolor{DefinitionBar}{HTML}{4ADE80}
\definecolor{ProblemBar}{HTML}{60A5FA}

\tcbset{
    blockbase/.style={
        enhanced,
        breakable,
        parbox=false,
        boxrule=0pt,
        left=10pt, right=10pt, top=10pt, bottom=10pt,
        arc=2pt,
    },
    leftbar/.style={
        borderline west={2.8pt}{0pt}{#1},
    },
    tinted/.style={
        colback=#1!18!PageBackground,
        colframe=#1!18!PageBackground,
        coltext=PageText,
        leftbar=#1,
    },
}

\makeatletter
\def\tagform@#1{\maketag@@@{[\ignorespaces#1\unskip\@@italiccorr]}}
\renewcommand{\eqref}[1]{\textup{[\ref{#1}]}}
\makeatother

\renewcommand{\phi}{\varphi}
\renewcommand{\epsilon}{\varepsilon}

\def\mbb#1{\mathbb{#1}}
\def\mfk#1{\mathfrak{#1}}
\def\bN{\mbb{N}}
\def\bC{\mbb{C}}
\def\bR{\mbb{R}}
\def\bQ{\mbb{Q}}
\def\bZ{\mbb{Z}}
\def\bF{\mbb{F}}
\def\bP{\mbb{P}}

\DeclarePairedDelimiter\floor{\lfloor}{\rfloor}
\DeclarePairedDelimiter\ceil{\lceil}{\rceil}
\DeclarePairedDelimiter{\angled}{\langle}{\rangle}
\let\norm\relax
\DeclarePairedDelimiter{\norm}{\lVert}{\rVert}

\DeclareMathOperator{\Span}{Span}
\DeclareMathOperator{\End}{End}
\DeclareMathOperator{\Hom}{Hom}
\DeclareMathOperator{\im}{im}
\DeclareMathOperator{\proj}{proj}
\DeclareMathOperator{\orth}{orth}
\DeclareMathOperator{\diag}{diag}

\newcommand{\rref}{\xrightarrow{\text{RREF}}}
\newcommand{\basis}{\mathcal}
\newcommand{\conj}{\overline}
\newcommand{\const}{\mathrm{const}}
\newcommand{\sym}{\mathrm{sym}}
\newcommand{\alt}{\mathrm{alt}}
\newcommand{\perm}{\mathrm{perm}}
\newcommand{\sgn}{\mathrm{sgn}}
\newcommand{\lcm}{\mathrm{lcm}}

\renewcommand{\div}{\divisionsymbol}
\renewcommand{\qedsymbol}{\textcolor{PageText}{$\square$}}

\newcommand{\innerli}{
    \par\vspace{-\parskip}\vspace{4pt}\noindent
    \makebox[\linewidth]{
        \textcolor{PageText!30}{\rule{0.47\linewidth}{0.5pt}}
        \hfill$\therefore$\hfill
        \textcolor{PageText!30}{\rule{0.47\linewidth}{0.5pt}}
    }
    \par\vspace{-\parskip}\vspace{4pt}
}

\newcommand{\ProofBodySpacing}{%
    \begingroup
    \setlength{\parskip}{1em}%
    \setlength{\parindent}{0pt}%
}
\newcommand{\EndProofBodySpacing}{\endgroup}

\newcommand{\tcbname}[1]{%
    \IfNoValueTF{#1}{}{%
        \if\relax\detokenize{#1}\relax\else #1\fi
    }%
}

\newcommand{\tcbheaderline}[2]{%
    \refstepcounter{equation}%
    {\color{#1}\textbf{[\theequation]}}\hspace{0.75em}%
    {\textsc{#2}}%
}

\newcommand{\tcbheadername}[1]{%
    \unskip
    \if\relax\detokenize{#1}\relax
    .\hspace{0.2em}%
    \else
        \hspace{0.4em}(\textit{#1}).\hspace{0.2em}%
    \fi
}

\NewDocumentEnvironment{axiom}{O{}}{%
    \begin{tcolorbox}[blockbase,tinted=AxiomBar]
    \tcbheaderline{AxiomBar}{Axiom}\tcbheadername{#1}\ignorespaces
}{%
    \unskip
    \end{tcolorbox}
}

\NewDocumentEnvironment{definition}{O{}}{%
    \begin{tcolorbox}[blockbase,tinted=DefinitionBar]
    \tcbheaderline{DefinitionBar}{Definition}\tcbheadername{#1}\ignorespaces
}{%
    \unskip
    \end{tcolorbox}
}

\NewDocumentEnvironment{remark}{O{}}{%
    \begin{tcolorbox}[blockbase,tinted=CorollaryBar]
    \tcbheaderline{CorollaryBar}{Remark}\tcbheadername{#1}\ignorespaces
}{%
    \unskip
    \end{tcolorbox}
}

\NewDocumentEnvironment{example}{O{}}{%
    \begin{tcolorbox}[blockbase,tinted=CorollaryBar]
    \tcbheaderline{CorollaryBar}{Example}\tcbheadername{#1}\ignorespaces
}{%
    \unskip
    \end{tcolorbox}
}

\NewDocumentEnvironment{theorem}{O{} +m}{%
    \begin{tcolorbox}[blockbase,tinted=TheoremBar]
    \tcbheaderline{TheoremBar}{Theorem}\tcbheadername{#1}{#2}
    \innerli
    \textbf{Proof.}\ \ignorespaces
    \ProofBodySpacing
}{%
    \unskip
    \EndProofBodySpacing
    \ifhmode
        \hfill\qedsymbol
    \else
        \vspace{-\baselineskip}\hfill\qedsymbol
    \fi
    \end{tcolorbox}
}

\NewDocumentEnvironment{proposition}{O{} +m}{%
    \begin{tcolorbox}[blockbase,tinted=PropositionBar]
    \tcbheaderline{PropositionBar}{Proposition}\tcbheadername{#1}{#2}
    \innerli
    \textbf{Proof.}\ \ignorespaces
    \ProofBodySpacing
}{%
    \unskip
    \EndProofBodySpacing
    \ifhmode
        \hfill\qedsymbol
    \else
        \vspace{-\baselineskip}\hfill\qedsymbol
    \fi
    \end{tcolorbox}
}

\NewDocumentEnvironment{lemma}{O{} +m}{%
    \begin{tcolorbox}[blockbase,tinted=LemmaBar]
    \tcbheaderline{LemmaBar}{Lemma}\tcbheadername{#1}{#2}
    \innerli
    \textbf{Proof.}\ \ignorespaces
    \ProofBodySpacing
}{%
    \unskip
    \EndProofBodySpacing
    \ifhmode
        \hfill\qedsymbol
    \else
        \vspace{-\baselineskip}\hfill\qedsymbol
    \fi
    \end{tcolorbox}
}

\NewDocumentEnvironment{lemma*}{O{} +m}{%
    \begin{tcolorbox}[blockbase,tinted=LemmaBar]
    \tcbheaderline{LemmaBar}{Lemma}\tcbheadername{#1}{#2}
    \end{tcolorbox}
}

\NewDocumentEnvironment{corollary}{O{} +m}{%
    \begin{tcolorbox}[blockbase,tinted=CorollaryBar]
    \tcbheaderline{CorollaryBar}{Corollary}\tcbheadername{#1}{#2}
    \innerli
    \textbf{Proof.}\ \ignorespaces
    \ProofBodySpacing
}{%
    \unskip
    \EndProofBodySpacing
    \ifhmode
        \hfill\qedsymbol
    \else
        \vspace{-\baselineskip}\hfill\qedsymbol
    \fi
    \end{tcolorbox}
}

\NewDocumentEnvironment{problem}{O{} +m}{%
    \begin{tcolorbox}[blockbase,tinted=ProblemBar]
    \tcbheaderline{ProblemBar}{Problem}\tcbheadername{#1}{#2}
    \innerli
    \textbf{Solution.}\ \ignorespaces
    \ProofBodySpacing
}{%
    \unskip
    \EndProofBodySpacing
    \ifhmode
        \hfill\qedsymbol
    \else
        \vspace{-\baselineskip}\hfill\qedsymbol
    \fi
    \end{tcolorbox}
}

\newcommand{\pset}[1]{
    \section*{\textbf{#1}}
    \addcontentsline{toc}{section}{#1}
}

\usepackage{chngcntr}
\numberwithin{equation}{subsection}

\ExplSyntaxOn
\NewDocumentCommand{\colvec}{m}
{
    \begin{bmatrix}
    \clist_use:nn { #1 } { \\ }
    \end{bmatrix}
}
\NewDocumentCommand{\rowvec}{m}
{
    \begin{bmatrix}
    \clist_use:nn { #1 } { & }
    \end{bmatrix}
}
\NewDocumentCommand{\blockmat}{m m}
{
    \left[\begin{array}{#1}
    #2
    \end{array}\right]
}
\ExplSyntaxOff

\newcommand{\splitlineh}[2]{%
    \rule{#1}{0.4pt}\,#2\,\rule{#1}{0.4pt}%
}

\newcommand{\splitlinev}[2]{%
    \begin{array}{c}
    \rule{0.4pt}{#1} \\[-2pt]
    #2 \\[-2pt]
    \rule{0.4pt}{#1}
    \end{array}
}

\AtBeginDocument{
    \setlength{\abovedisplayskip}{6pt plus 1pt minus 1pt}
    \setlength{\belowdisplayskip}{6pt plus 1pt minus 1pt}
    \setlength{\abovedisplayshortskip}{0pt plus 1pt}
    \setlength{\belowdisplayshortskip}{4pt plus 1pt}
    \setlength{\jot}{4pt}
    \allowdisplaybreaks
    \setlength{\parindent}{0pt}
    \setlength{\parskip}{1em}
}

\usepackage{hyperref}
\hypersetup{
    colorlinks=true,
    linkcolor=ProblemBar,
    filecolor=ProblemBar,
    urlcolor=ProblemBar,
    citecolor=ProblemBar
}
